\section{Threats to Validity}
\label{sec:threats}

\paragraph{Two benchmarks, one natural procedure.} The contract's form adds
a few points only on tasks governed by a rule-matching procedure, and the
only natural one in either benchmark is DABStep's fee matching. Experiment
4 builds three more, but they are ours, and complete prose carries them at
the ceiling. A second, independent benchmark whose documentation states
rules once for many tasks would test the form result properly; EntSQL,
DataSpace and BI-Bench give nearly every task its own document or project,
leaving no rules to state once~\cite{entsql,dataspace,bi-bench}.

\paragraph{The raters are models.} The raters who judged determinacy,
procedural load and the gaps of experiment 4b, and the failure auditors,
are instances of one model family. Each analysis plan was fixed
before any rating was read, and the raters saw no run, contract or result,
but their agreement ($\kappa$ 0.68 to 0.96) overstates what human raters
would reach; on one batch two raters matched on all 71 loads with
different notes.

\paragraph{Ceilings.} Experiments 4 and 6 sit at the ceiling at their base
level, so they could show a drop but not a gain, and their constructed
questions are easier than the benchmarks' own. Experiment 6's padding is
other databases under distinct names; neighbouring domains of one warehouse
would be the harder test of a long prompt.

\paragraph{Repetition.} \armTws{} states each fact twice, once in the
contract and once in the note, against the prose's once, so part of the
form gain of Section~\ref{sec:alike} may be repetition.

\paragraph{Scoring views and graders.}\label{sec:threats-golds} \mQwen{} states its working before
its answer, which the official DABStep scorer marks wrong; we report the
end-to-end view beside it, and on experiment 5b the order of the arms
changes between the two views. The vendored DABStep scorer also ignores the
sign of a number: on nine fee-change questions with negative golds, 114
official verdicts and 43 end-to-end ones credit an answer of the wrong
sign, mostly in \mQwen{}'s baseline arms. Re-grading with the sign enforced
changes no prediction's outcome; it moves \mQwen{}'s scaffolding step from
$-4.3$ to $-3.4$ points (official) and from $+0.2$ to $+1.2$
(end-to-end), and we keep the official verdicts. LiveSQLBench's official
rules strip every \texttt{DISTINCT} from both answer and gold, which leans
against the compiled contracts' \texttt{IS DISTINCT FROM}; grading without
the stripping moves each arm by under two points and no outcome. DABStep's
reconstructed golds run 3 to 6 points high on the hard split, on both
arms of the officially graded pair.

\paragraph{Variance and multiplicity.}\label{sec:power} A near-replicate of the \mGLM{}
contract arm at temperature~0 disagreed with the original on 94 of 401
tasks, symmetrically; the four single runs are therefore read only through
paired tests, and the repeated experiments carry the claims. The GPT-6
models accept only their default temperature. No multiplicity correction
is applied across the 53 predictions; that 20 of them failed is the
clearest sign they were not fitted.

\paragraph{Contamination.} Both benchmarks are public, and we ran no
contamination probe. A memorising model should not care which context it
is handed, yet the hollow and full contracts, which differ only in their
prose, differ by 26 to 50 points on the same tasks and model. That is an inference, not a measurement.

\paragraph{Run health.} \mDS{}'s single run lost 29\% of its rows to rate
limiting and is reported on its complete-case set. In experiment 6, 174 to
180 of each arm-level's 180 runs kept a trace; the rest ran before a pod
restart. \mSol{}'s text arm at 8x and all 22 databases was dropped for
cost by a recorded amendment, before any test.
